\honest{Rationalization}{Eliezer Yudkowsky}{2007}

I recap ``The Bottom Line,'' from two days ago. The curious inquirer lists the clues and then concludes; the clever arguer concludes and then picks the clues. In between, yesterday, in ``What Evidence Filtered Evidence?'', I said that each of the clever arguer's statements is valid evidence. Here I again treat the arguer as the pure opposite of reason.

The first procedure, I say, is rationality. The second is called ``rationalization,'' and I object to the word. ``You cannot `rationalize' what is not already rational,'' I say; it is like calling lying ``truthization.'' But to rationalize something means to make it seem reasonable, which is exactly what the arguer does. I ask, ``What fool devised such confusingly similar words,'' and propose ``giant sucking cognitive black hole'' instead. The word gets two paragraphs.

I restate the difference as a numbered list. Either you start from the evidence and work forward to a conclusion, or you start from a conclusion and work backward to pick the evidence.

Rationality, I say, changes beliefs to make them more accurate, while rationalization fixes them in place, so it should be called ``anti-rationality.'' But an honest check can also end where it began. What matters is whether the belief could have moved, not whether it did.

I restate the forward direction, adding that in probability theory your expected future belief equals your present one: ``If you know your destination, you are already there.'' Then I restate the backward direction.

Then I blame ``Traditional Rationality,'' which, I say, approves of a scientist who picks a pet hypothesis and sets out to prove it. I name no one who holds this view. The approving line is spoken by a Traditional Rationalist I made up. Karl Popper, the best-known philosopher of scientific method, said that a genuine test of a theory is an attempt to refute it. Then I concede the point anyway: pride ``is the engine that drives Science forward,'' because it is easier to find two people biased in opposite directions than one person who is unbiased.

My reply is that ``just because everyone does something, doesn't make it okay.'' Nobody argued that everyone does it. The argument was that it works. It would be better, I say, if the scientist set out to test the hypothesis out of curiosity. But that scientist still starts from the conclusion, which by my own definition is working backward. What makes the scientist rational is the test. So the direction of the reasoning was never what mattered.

If you do not know where you are going, I say, you will probably feel curious about it. People are ignorant of many things they do not care about.

I end: ``Feel the flow of the Force, and make sure it isn't flowing backwards.'' I give no way to check which way it is flowing.
